% =============================================================================
% cap3_pipeline/03_fase2_deteccao.tex
% Seção 3.3 -- Fase II: Detecção generalista de animais
%   Modelo: MegaDetector (pré-treinado, as-is), classe agnóstica.
%   Função: descartar agressivamente quadros sem animal antes da
%   classificação multi-espécie da Fase III.
%   Texto corrido, impessoal, sem listas no corpo.
% =============================================================================
\section{Fase II --- Detecção generalista}
\label{sec:pipeline-fase2}

A Fase~II decide, para cada quadro canônico recebido da Fase~I, se há ou não um animal presente e, em caso afirmativo, localiza a sua \emph{bounding box}. Opera como uma válvula entre uma entrada potencialmente abundante (cartões com milhares de quadros amostrados) e uma saída relativamente escassa (quadros que efetivamente contêm um animal junto ao ponto de alimentação). Em projetos de \emph{camera trapping} em larga escala é convergente o achado de que a maior parte do volume bruto coletado consiste em \emph{ghost frames} acionados por vento, sombra ou variação térmica, e que filtrar esse volume antes de qualquer etapa de classificação fina é condição para a viabilidade operacional do restante do pipeline~\cite{bruce2025largescale,velez2022choosing,swanson2015snapshot,mulero2025addressing}.

\subsection*{Escolha do modelo: \modeloMD~como detector base}

O modelo adotado nesta fase é o \modeloMD, integrado por meio da pilha \emph{PyTorch-Wildlife} mantida pelo grupo \emph{AI for Wildlife} da Microsoft Research~\cite{hernandez2024pytorchwildlife,microsoft_cameratraps_repo}. Trata-se de um detector treinado em grande volume de dados de \emph{camera trapping} de diferentes biomas, projetado para detecção generalista de fauna em três classes amplas (animal, pessoa, veículo). Sua adoção é justificada por três fatores convergentes: disponibilidade pública dos pesos e da \emph{pipeline} de inferência, maturidade documentada em campo e alinhamento natural com a Fase~III, que assume para si a classificação multi-espécie. Esse desenho evita misturar em um único modelo o problema de rejeição de vazio e o de classificação fina, com modos de falha distintos e avaliação separável.

Em coerência com a decisão de escopo da Seção~\ref{sec:pipeline-visao-geral}, o \modeloMD~é utilizado sem \emph{fine-tuning} sobre dados do Campus~2: o detector já cobre, com folga, a classe macroscópica ``animal'' que interessa nesta fase, e o esforço de coleta e anotação de um dataset local específico para detecção não traria ganho proporcional ao custo no escopo deste trabalho. Eventuais adaptações de domínio --- por exemplo, ajuste sobre quadros NIR noturnos, regime em que o \modeloMD~tem desempenho documentadamente inferior ao diurno --- ficam registradas como trabalhos futuros no Capítulo~\ref{cap:conclusao}.

\subsection*{Operação da fase no pipeline}

Para cada quadro canônico entregue pela Fase~I, o \modeloMD~produz uma lista de detecções, cada uma composta por uma \emph{bounding box}, uma classe macroscópica e um \emph{score} de confiança. A Fase~II aplica sobre essa saída um limiar de confiança configurável: quadros sem detecção acima do limiar são considerados negativos e removidos do fluxo; quadros com ao menos uma detecção válida da classe ``animal'' são marcados como positivos e seguem para a Fase~III, acompanhados das coordenadas da \emph{bounding box}, o que permite recortar a região de interesse e poupar à Fase~III o custo de processar o quadro inteiro. Detecções das classes ``pessoa'' e ``veículo'', quando presentes, são registradas apenas como metadado de diagnóstico (causa provável de disparo) e não propagam como entrada para a Fase~III.

\subsection*{Anonimização de pessoas na origem}

A presença incidental de transeuntes nos quadros é tratada por um passo de mascaramento automático aplicado imediatamente após a inferência do \modeloMD~e antes de qualquer persistência em disco: as regiões correspondentes a detecções da classe ``pessoa'' são submetidas a um filtro gaussiano de alta intensidade, de modo que apenas versões anonimizadas dos quadros sejam gravadas no acervo do projeto. As coordenadas da \emph{bounding box} permanecem como metadado operacional, sem qualquer associação a indivíduos identificáveis. O arranjo é coerente com o princípio de \emph{privacy-by-design} adotado pelo projeto e atende, na origem, às obrigações decorrentes da Lei Geral de Proteção de Dados (Lei n\textsuperscript{o}~13.709/2018); o detalhamento jurídico-operacional foge ao escopo técnico deste capítulo.

\subsection*{Limites conhecidos da Fase II e mitigações}

A literatura aponta dois limites principais do \modeloMD~em configuração padrão. O primeiro é o desempenho relativamente inferior em quadros noturnos sob iluminação NIR, em que o contraste entre animal e fundo é menor e o detector tende a produzir mais falsos negativos~\cite{velasco2024smartcameratraps,mulero2025addressing}; a mitigação adotada é configurar o limiar de confiança por modo de captura (diurno versus noturno) e estratificar a avaliação do Capítulo~\ref{cap:metodologia} por horário, para que eventuais quedas de \mAP~em regime noturno apareçam explicitamente nos resultados. O segundo é o comportamento em \emph{full-frame versus crop}: quando o animal ocupa fração muito pequena ou muito grande da imagem, o detector pode produzir caixas imprecisas ou falhar; a mitigação é parametrizar tamanhos mínimo e máximo de \emph{bounding box} e estratificar a avaliação por distância média ao animal por ponto, sem alteração do modelo em si. Ao final da Fase~II, o que segue para a Fase~III é um subconjunto significativamente menor de quadros, cada um acompanhado da região de interesse sobre a qual a classificação multi-espécie precisa operar.
